\section{Experimental Method and Evidence}
\subsection{Frozen Inputs, Cohorts and Resources}
Construction-only controls use separately bound RRGs; capped route tails do not establish legality. Historical native BNN/s38417 compare Stock, Flat, ALT CPU and original ALT CUDA. LSTM seed2 adds Pipeline on another cohort. Previously used in the CPU study, LSTM is a same-family holdout from preceding ALT tuning, frozen before its own tests, not before all ALT development. It shares BNN's architecture but has a different netlist, placement and RRG.

New matched GEMM/BNN runs retain their original seed1 inputs: width100 for GEMM, width120 for BNN. LSTM also uses width120; all three have two layers and 40 samples (Table~\ref{tab:qor}). The new allocation fixes eight physical cores and NUMA0 memory binding on a Xeon Platinum 8462Y+, one H100 80GB, a 10-CPU quota and 200GiB RAM. The earlier native ALT cohorts use eight physical cores in one NUMA node on Gold 6530 and one H100 80GB. Solvers execute sequentially, with rotated arm order and fresh processes. Shared hosts are not exclusive; cohorts are not pooled. The baseline is VTR commit \texttt{0b912334fdd8}, release build, assertion level 2.

\subsection{Timing and Statistical Policy}
E1 is sample service; CUDA ends it before final ABI output copy. E2 is native \texttt{MapLookahead::compute}: conversion, preparation/service, output handoff, group merge, interpolation, source/OPIN and derived minima. Process adds RR XML loading, initialization, routing, checking/STA and uniform serialization. Independent route readback, pre-generated RRG, packing and placement are excluded. Caches are not flushed.

For paired repeat $r$, $s_r=T_{c,r}/T_{a,r}$. We report mean and sample SD of ratios and raw seconds ($n=3$). SD is not a confidence interval; repeats are processes, not designs or seeds. Stock E2 precision is 0.01s; other E2 timers are high resolution. Profiles are outside formal timing.

\subsection{Qualification and Attribution}
STRICT checks wire-map bytes under its proposition. ALT checks complete CPU/CUDA pairs and maps, then each route through independent original-VPR FULL/STA with unchanged route SHA, zero overuse, and CPD/wirelength within $+5\%$ of controls: Flat/Stock for LSTM, Stock/Pipeline for new GEMM/BNN. New ALT CPU runs are qualification-only; formal CUDA disables capture and CPU-reference verification.

Retained bounding-box notification repairs can change route trajectories despite identical maps; a unique complete cause is not established. Stock ratios describe the implementation combination. Shared-repair Flat/Pipeline/ALT comparisons better isolate construction.
